% !TeX root = ../PIFB2.tex
\section{From the Statistical Theory to a Physical Program}\label{sec:part2_preamble}

Part~I supplies a normalized population model, selective communication, variational inference and conditional constructions of retained agents. Part~II asks how structures of those types could acquire physical meaning. The question is stronger than whether a familiar physical equation can be written using information-theoretic symbols. A physical interpretation must connect a state of the model to a reproducible record, derive how that record responds to interventions, and identify which alternative models would give a different result. Wheeler's participatory program motivates this direction~\cite{Wheeler1990}; the mathematics does not make its philosophical interpretation compulsory.

Three distinctions organize the development. A mathematical construction specifies an object and proves its properties. Fisher length is a length of a statistical history, the covariant Fisher pullback is a positive semidefinite base tensor, and the Hermitian determinant construction carries the full proper Lorentz group. An operational attachment specifies how such an object is read from records, including the measurement channel and calibration. A dynamical selection result establishes that the updating law produces or stabilizes the object under stated initial conditions. These are separate obligations. A tensor with Lorentzian signature meets the first; it becomes a spacetime candidate only after a nondegenerate collective readout and a compatible propagation law are supplied.

The existing results already reduce the number of independent constructions required. Parent laws arise by retaining variables of a fine law; upward and downward cross-scale operators follow from pushforward and conditional reconstruction; priors for a later episode follow from prediction through a declared transition kernel. Effective interactions arise from integrating out fine variables, and conditional memory records what an instantaneous parent law omits. None of these operations needs a further shadow force or an additional substance. Conversely, exact marginalization alone does not explain why a particular parent forms, why it persists, or why the resulting observable has a physical unit. Those are the next questions, not definitions to be inserted into the answer.

The same economy guides the geometric discussion. The one statistical bundle of Section~\ref{sec:statistical_manifolds} is sufficient for joint, marginal and predictive information geometry. A connection is declared only when comparing section values across the continuous base. A complex rank-two channel is introduced explicitly when deriving the Lorentz representation, and a soldering is required only when attaching that representation to base tangents. We do not add an indefinite frame-derivative block to the Fisher metric or infer a physical dimension count from a threshold chosen for that purpose. Each specialization is stated where it is used rather than collected into a second hierarchy of universal postulates.

The companion theory provides exact identities, regularity conditions and obstructions for these constructions~\cite{DennisGaugeVFE2026,DennisGaugeVFESuppB2026}. Its mathematical scope neither proves nor refutes a proposed physical realization. The task here is to carry its usable structure into a physicist's account while retaining the distinction between a proved conditional result and a physical hypothesis. The following sections develop duration, observer and collective geometry, Lorentz kinematics, persistence, measurement and empirical symmetry in that order. Consciousness is discussed separately as an interpretation on which none of those derivations depends.
